%% ChurnBench: A Drift-Aware Benchmark for Grounding Agentic AI
%%             over Enterprise Data Fabrics
%%
%% IEEE Access submission draft.
%%
%% PLACEHOLDER CONVENTION
%%   [DATA REQUIRED]    — a number or table that must come from a committed
%%                        experiment artifact in results/*.json before submission.
%%                        NEVER fill with invented or estimated figures.
%%   [CITATION NEEDED]  — a claim that needs a real reference; add to
%%                        references.bib and replace this marker.
%%   [FIGURE PENDING]   — figure exists in paper/figures/ or is described in
%%                        a TODO comment; must be finalised before submission.
%%
%% Build: latexmk -pdf churnbench-paper.tex
%% Requires: ieeeaccess.cls (copy from IEEE Access author kit — not committed)
%%           IEEEtran.bst   (LPPL — committed)
%%
\documentclass{ieeeaccess}

\usepackage{cite}
\usepackage{amsmath,amssymb,amsfonts}
\usepackage{algorithmic}
\usepackage{graphicx}
\usepackage{textcomp}
\usepackage{xcolor}
\usepackage{booktabs}
\usepackage{multirow}
\usepackage{tabularx}
\usepackage{listings}
\usepackage{hyperref}

\lstset{
  basicstyle=\ttfamily\footnotesize,
  breaklines=true,
  columns=fullflexible,
  keepspaces=true,
  frame=single,
  framesep=4pt,
  xleftmargin=4pt,
  showstringspaces=false,
}

\graphicspath{{figures/}}

% ── Title & Authors ───────────────────────────────────────────────────────────
\title{ChurnBench: A Drift-Aware Benchmark for Grounding Agentic AI
       over Enterprise Data Fabrics}

\author{%
  \uppercase{Vivek Kumar Singh}\authorrefmark{1}
  \IEEEmembership{Independent Researcher}
}

\address[1]{McKinney, TX, USA (e-mail: vivekksingh.nov12@gmail.com)}

% ── Abstract ─────────────────────────────────────────────────────────────────
\begin{abstract}
Agentic AI systems that answer questions over enterprise data must
reason across heterogeneous, continuously mutating sources — relational
warehouses, operational document stores, rate-limited SaaS APIs, and
unstructured document corpora.  Existing benchmarks evaluate static
snapshots, making it impossible to measure \emph{freshness error}: the
gap between an agent's answer at query time $T$ and the true world
state at $T$.

We introduce \textbf{ChurnBench}, an open-source evaluation instrument
for this setting.  ChurnBench simulates a Software Asset Management
(SAM) enterprise fabric over a configurable time horizon, writing every
mutation to an append-only ground-truth \emph{ledger}.  The ledger can
be projected onto live fabric stores — Postgres warehouse, MongoDB
operational store, a rate-limited FastAPI mock-SaaS, and a synthetic
contract document corpus — at any frozen timestamp $T$, enabling
deterministic, reproducible evaluation.

We define a freshness error metric, describe a suite of
[DATA REQUIRED]~benchmark tasks spanning point lookups, aggregations,
and multi-hop joins across fabric layers, and evaluate
[DATA REQUIRED]~agent architectures.  The best-performing arm achieves
[DATA REQUIRED]\% freshness accuracy at $T$ while reducing SaaS API
calls by [DATA REQUIRED]$\times$ relative to a naive tool-calling
baseline.

ChurnBench is available at \url{https://github.com/vsingh45/churnbench}.
\end{abstract}

\begin{keywords}
agentic AI, grounding, enterprise data fabric, temporal drift,
benchmark, software asset management, freshness error, RAG
\end{keywords}

\titlepgskip=-15pt
\maketitle

% ── 1. Introduction ──────────────────────────────────────────────────────────
\section{Introduction}
\label{sec:intro}

Enterprise AI agents face a distinctive challenge: the data they must
reason over is not static.  Software licenses get reassigned;
contracts renew; users offboard; prices change.  An agent that answers
``Which users hold an active license for Product~X?'' correctly at
$T = \text{Monday}$ may give a stale answer by $T = \text{Wednesday}$
if a licence was reassigned on Tuesday.

We call this \emph{temporal drift} and the resulting evaluation error
\emph{freshness error}.  Freshness error is systematically unmeasured
by existing benchmarks [CITATION NEEDED] because those benchmarks evaluate
agents against a static database snapshot — they provide no mechanism
to vary $T$ and compare the agent's answer against the \emph{true}
world state at $T$.

This paper makes the following contributions:

\begin{itemize}
  \item \textbf{ChurnBench}, an open-source, fully deterministic
    benchmark generator for grounding agentic AI over heterogeneous,
    drifting enterprise data (Section~\ref{sec:benchmark}).

  \item A formal definition of \emph{freshness error} and a measurement
    methodology based on the ground-truth ledger
    (Section~\ref{sec:metric}).

  \item A suite of [DATA REQUIRED]~benchmark tasks across four
    difficulty tiers and three fabric layers
    (Section~\ref{sec:tasks}).

  \item Evaluation of [DATA REQUIRED]~agent architectures on
    ChurnBench, with analysis of the accuracy--cost tradeoff under
    rate-limited API access (Section~\ref{sec:experiments}).
\end{itemize}

% ── 2. Related Work ───────────────────────────────────────────────────────────
\section{Related Work}
\label{sec:related}

\subsection{Enterprise QA and Text-to-SQL Benchmarks}
Spider~[CITATION NEEDED], BIRD~[CITATION NEEDED], and
WikiSQL~[CITATION NEEDED] evaluate structured-query generation over
static relational databases.  They do not model temporal mutation,
multiple storage tiers, or rate-limited API access.

\subsection{Retrieval-Augmented Generation}
RAG benchmarks such as BEIR~[CITATION NEEDED] and RAGAS~[CITATION NEEDED]
evaluate retrieval quality over document corpora.  ChurnBench extends
this to multi-modal enterprise fabrics where documents, structured
records, and API responses must be jointly reasoned over.

\subsection{Agent Benchmarks}
WebArena~[CITATION NEEDED], AgentBench~[CITATION NEEDED], and
GAIA~[CITATION NEEDED] evaluate autonomous agents on web and tool-use
tasks.  None model the setting where the ground truth itself drifts
over time.

\subsection{Temporal Reasoning}
TempReason~[CITATION NEEDED] and TimeQA~[CITATION NEEDED] probe
temporal reasoning in language models but use static datasets; they
do not evaluate live grounding under mutation.

% ── 3. The ChurnBench Benchmark ───────────────────────────────────────────────
\section{The ChurnBench Benchmark}
\label{sec:benchmark}

\subsection{Domain: Software Asset Management}
\label{sec:domain}

SAM is a natural testbed for drift-aware grounding: licenses are
continuously reassigned, contracts renew on staggered cycles, users
onboard and offboard daily, and prices change on vendor schedules.
The temporal structure is well-defined and the ground truth is
binary (a user either holds a license at time $T$ or does not),
making scoring straightforward.

\subsection{The Data Fabric}
\label{sec:fabric}

ChurnBench materialises four storage tiers, mirroring common
enterprise architectures:

\begin{description}
  \item[Postgres 16 warehouse]  A star-schema-ish analytical store
    (\texttt{sam.*}) with dimension tables for vendors, products, cost
    centers, and fact tables for license purchases and consumption
    events.  Historical; append-mostly.

  \item[MongoDB 7 operational store]  Current state collections
    (\texttt{sam\_ops.*}): \texttt{users}, \texttt{active\_licenses},
    \texttt{assignments}, \texttt{entitlements}, \texttt{tickets},
    \texttt{utilization\_current}.  Mutable; reflects live business
    state.

  \item[Mock SaaS API]  A FastAPI service with configurable latency
    ($350\,\text{ms}$ default) and rate limit ($60\,\text{req/min}$
    default), modelling the cost of hitting a real third-party
    integration per query.

  \item[Document corpus]  Synthetic contract markdown files, one per
    active contract at time $T$.  Renewal events update the effective
    term and expiry fields.
\end{description}

\subsection{The Ground-Truth Ledger}
\label{sec:ledger}

All fabric mutations are generated by a discrete-day simulator
(\texttt{TimelineSimulator}) and written to an append-only
\emph{ledger} of typed events (Table~\ref{tab:events}).  The ledger
is the sole source of truth: evaluating an agent's answer at $T$
means comparing it to the state derived by folding
\texttt{events\_through(T)}.

\begin{table}[t]
  \caption{Ledger Event Types}
  \label{tab:events}
  \centering
  \begin{tabular}{ll}
    \toprule
    \textbf{Event Kind} & \textbf{Effect} \\
    \midrule
    \texttt{USER\_HIRED}              & Add user to active set \\
    \texttt{USER\_OFFBOARDED}         & Remove user; unassign licenses \\
    \texttt{USER\_MOVED\_COST\_CENTER} & Update user's cost center \\
    \texttt{LICENSE\_PURCHASED}        & Add license record \\
    \texttt{LICENSE\_ASSIGNED}         & Set holder \\
    \texttt{LICENSE\_UNASSIGNED}       & Clear holder \\
    \texttt{LICENSE\_REASSIGNED}       & Swap holder \\
    \texttt{CONTRACT\_SIGNED}          & Create contract record \\
    \texttt{CONTRACT\_RENEWED}         & Update term / effective date \\
    \texttt{PRICE\_CHANGED}            & Update unit price \\
    \texttt{CONSUMPTION\_LOGGED}       & Append consumption event \\
    \bottomrule
  \end{tabular}
\end{table}

\subsection{Determinism and Reproducibility}
\label{sec:determinism}

Given a fixed \texttt{(seed, config)}, the simulator produces an
identical ledger on any machine.  The projector rebuilds each fabric
layer by wiping and replaying events through $T$ — no partial updates.
Two runs at the same $T$ produce bit-identical stores.

% ── 4. Freshness Error Metric ─────────────────────────────────────────────────
\section{Freshness Error Metric}
\label{sec:metric}

Let $q$ be a benchmark task with ground-truth answer $a^*(T)$
derived from \texttt{ledger.events\_through(T)}.  Let $\hat{a}(T)$ be
the agent's answer at query time $T$.

\begin{definition}[Freshness Error]
  $\text{FE}(q, T) = \mathbf{1}[\hat{a}(T) \neq a^*(T)]$
\end{definition}

\begin{definition}[Freshness Accuracy]
  $\text{FA} = 1 - \frac{1}{|Q||T|}\sum_{q \in Q}\sum_{T \in \mathcal{T}}
  \text{FE}(q, T)$
\end{definition}

where $Q$ is the task suite and $\mathcal{T}$ is a set of evaluation
timestamps spanning the simulation horizon.

% ── 5. Task Suite ────────────────────────────────────────────────────────────
\section{Task Suite}
\label{sec:tasks}

% [FIGURE PENDING] — task taxonomy diagram (paper/figures/fig_tasks.pdf)

Tasks are generated from the ledger at evaluation time and span four
difficulty tiers:

\begin{description}
  \item[T1 — Point lookup]  Single-entity, single-fabric queries
    (``Is usr\_000042 active?'', ``What is the current price of
    prd\_0007?'').  Ground truth is a single value from one ledger
    fold.

  \item[T2 — Filtered aggregate]  Count or sum over a filtered set
    (``How many active licenses does cost center cc\_003 hold?'').
    Requires joining across two Mongo collections or a Postgres fact
    table.

  \item[T3 — Cross-fabric join]  Questions that require synthesising
    data from at least two fabric tiers (``Which users hold a license
    for a product whose price changed in the last 30 days, and what
    are their current entitlements?'').

  \item[T4 — Temporal delta]  Questions about what changed between
    $T_1$ and $T_2$ (``Which licenses were reassigned between
    2024-01-15 and 2024-02-01?'').  Requires
    \texttt{ledger.events\_between(T1, T2)}.
\end{description}

The benchmark will include [DATA REQUIRED]~tasks total:
[DATA REQUIRED]~T1, [DATA REQUIRED]~T2, [DATA REQUIRED]~T3,
[DATA REQUIRED]~T4.

% ── 6. Agent Architectures ───────────────────────────────────────────────────
\section{Agent Architectures Under Evaluation}
\label{sec:arms}

We evaluate [DATA REQUIRED]~arms:

\begin{description}
  \item[Arm 1 — Naive tool-calling]  Direct tool calls to the live
    fabric for every sub-question.  Worst-case freshness accuracy
    (always reads ``now'', not at $T$) but maximum API cost.

  \item[Arm 2 — Snapshot-grounded]  Agent receives a frozen projection
    of the fabric at $T$ via the ChurnBench projector before
    answering.  Best-case freshness accuracy; no live API calls needed.

  \item[Arm 3 — Hybrid caching]  Agent reads from a cache warmed at
    the last projection $T_\text{prev}$, with selective live calls for
    entities flagged as likely-drifted.  Tradeoff arm.
    % [CITATION NEEDED] — cite related cache-invalidation / staleness
    % detection work here.
\end{description}

% ── 7. Experiments ───────────────────────────────────────────────────────────
\section{Experiments}
\label{sec:experiments}

% [DATA REQUIRED] — all numbers in this section come from
% results/freshness_accuracy.json and results/api_cost.json.
% Do not fill any value here until those files are committed.

\subsection{Setup}
[DATA REQUIRED] — hardware, model (claude-opus-4-8 / claude-sonnet-5),
simulation config (days, seed, n\_tasks), evaluation timestamps.

\subsection{Freshness Accuracy}

\begin{table}[t]
  \caption{Freshness Accuracy by Arm and Task Tier}
  \label{tab:results}
  \centering
  \begin{tabular}{lcccc}
    \toprule
    \textbf{Arm} & \textbf{T1} & \textbf{T2} & \textbf{T3} & \textbf{T4} \\
    \midrule
    Arm 1 (Naive)    & [DATA REQUIRED] & [DATA REQUIRED] & [DATA REQUIRED] & [DATA REQUIRED] \\
    Arm 2 (Snapshot) & [DATA REQUIRED] & [DATA REQUIRED] & [DATA REQUIRED] & [DATA REQUIRED] \\
    Arm 3 (Hybrid)   & [DATA REQUIRED] & [DATA REQUIRED] & [DATA REQUIRED] & [DATA REQUIRED] \\
    \bottomrule
  \end{tabular}
\end{table}

\subsection{API Cost Under Rate Limiting}

% [FIGURE PENDING] — paper/figures/fig_api_cost.pdf
% Results source: results/api_cost.json

[DATA REQUIRED] — API call counts and latency per arm.

\subsection{Error Analysis}

[DATA REQUIRED] — breakdown of freshness errors by event kind
(e.g.\ LICENSE\_REASSIGNED causes most T1 errors; CONTRACT\_RENEWED
drives most doc-layer errors).

% ── 8. Discussion ────────────────────────────────────────────────────────────
\section{Discussion}
\label{sec:discussion}

\subsection{Limitations}

ChurnBench currently models a single enterprise domain (SAM).
Generalisation to other drifting enterprise settings (ERP, CRM,
ITSM) is left to future work.  The mock SaaS API is intentionally
simplified; a real SaaS integration has authentication, pagination,
and webhook-push patterns not modelled here.

The document corpus uses synthetic contract text; real contract
language is more varied and ambiguous, which may inflate T3/T4
difficulty for document-layer sub-questions.

\subsection{Benchmark Validity}

The ledger-as-ground-truth design ensures that freshness error is
measured against the \emph{same} facts that generated the fabric
state, not against a separately-maintained oracle.  This closes a
common evaluation gap where oracle and data-source can disagree due
to synchronisation lag.

% ── 9. Conclusion ────────────────────────────────────────────────────────────
\section{Conclusion}
\label{sec:conclusion}

We have presented ChurnBench, a drift-aware benchmark for evaluating
agentic AI grounding over heterogeneous enterprise data fabrics.  The
ledger-projection architecture enables reproducible frozen-at-$T$
evaluation and makes freshness error a first-class, measurable quantity.

[DATA REQUIRED] — summary of key experimental findings once results
are available.

The benchmark and all evaluation code are open-source at
\url{https://github.com/vsingh45/churnbench}.

% ── Acknowledgements ─────────────────────────────────────────────────────────
\section*{Acknowledgements}
[CITATION NEEDED] — acknowledge any collaborators, compute resources,
or funding sources before submission.

% ── References ───────────────────────────────────────────────────────────────
\bibliographystyle{IEEEtran}
\bibliography{references}

\end{document}
